\documentclass[10pt,a4paper]{article}

% ============================================================
%  Jilin University — Stylized Technical Report Template
%  Compact two-column layout inspired by Google DeepMind papers
% ============================================================

% ---------- geometry (tight margins like NeurIPS/ICML) ----------
\usepackage[
  top=0.75in,
  bottom=0.85in,
  left=0.62in,
  right=0.62in
]{geometry}

\usepackage[utf8]{inputenc}
\usepackage[T1]{fontenc}
\usepackage{CJKutf8}

% ---------- fonts ----------
\usepackage{XCharter}
\usepackage[scaled=1.1]{zlmtt}
\usepackage{microtype}
\usepackage{helvet}

% ---------- math ----------
\usepackage{amsmath,amssymb,amsfonts}
\usepackage{mathtools}
\usepackage{bm}

% ---------- graphics & tables ----------
\usepackage{graphicx}
\usepackage{booktabs}
\usepackage{multirow}
\usepackage{subcaption}
\usepackage{wrapfig}
\usepackage{float}

% ---------- two-column ----------
\usepackage{multicol}
\setlength{\columnsep}{0.25in}

% ---------- algorithms ----------
\usepackage{algorithm}
\usepackage{algorithmic}

% ---------- spacing ----------
\usepackage{setspace}
\setstretch{0.95}
\setlength{\parskip}{2pt plus 1pt minus 0.5pt}
\setlength{\parindent}{1em}
\setlength{\textfloatsep}{8pt plus 2pt minus 2pt}
\setlength{\floatsep}{6pt plus 2pt minus 2pt}
\setlength{\intextsep}{6pt plus 2pt minus 2pt}
\setlength{\abovecaptionskip}{4pt}
\setlength{\belowcaptionskip}{2pt}

% ---------- compact lists ----------
\usepackage{enumitem}
\setlist{nosep,leftmargin=*,topsep=2pt,parsep=1pt,itemsep=1pt}

% ---------- hyperlinks & references ----------
\usepackage[table]{xcolor}

\definecolor{jluDeepBlue}{RGB}{0,51,153}
\definecolor{jluRed}{RGB}{198,40,40}
\definecolor{jluTeal}{RGB}{0,137,123}
\definecolor{jluAmber}{RGB}{255,160,0}
\definecolor{jluPurple}{RGB}{106,27,154}
\definecolor{jluLightBg}{RGB}{245,247,252}
\definecolor{jluGray}{RGB}{100,100,100}

\newcommand{\MonthYear}{\the\year-\ifnum\month<10 0\fi\the\month}

\usepackage[colorlinks=true,linkcolor=jluDeepBlue,citecolor=jluTeal,urlcolor=jluRed]{hyperref}
\usepackage{cleveref}
\usepackage[numbers,sort&compress]{natbib}
\bibliographystyle{unsrtnat}

% ---------- misc ----------
\usepackage{fancyhdr}
\usepackage{tikz}
\usetikzlibrary{calc}
\usepackage{tcolorbox}
\tcbuselibrary{skins,breakable}
\usepackage{titlesec}
\usepackage{caption}
\captionsetup{font=small,labelfont=bf,skip=3pt}

% ============================================================
%  PAGE STYLE
% ============================================================
\pagestyle{fancy}
\fancyhf{}
\renewcommand{\headrulewidth}{0.8pt}
\renewcommand{\headrule}{%
  \begingroup\color{jluDeepBlue}\hrule width\headwidth height 0.8pt\endgroup}
\fancyhead[L]{%
  \raisebox{-0.18ex}{\logoimg[2.05em]{university}{jluDeepBlue}{University}}}
\fancyhead[R]{\small\rmfamily\color{black}\MonthYear}
\fancyfoot[C]{\footnotesize\sffamily\color{jluGray}\thepage}

% ============================================================
%  SECTION HEADING STYLES — compact & colourful
% ============================================================
\titleformat{\section}
  {\normalsize\rmfamily\bfseries\color{jluDeepBlue}}
  {\colorbox{jluDeepBlue}{\textcolor{white}{\;\thesection\;}}}{0.5em}{}
  [\vspace{1pt}{\color{jluDeepBlue}\titlerule[0.8pt]}]

\titleformat{\subsection}
  {\normalsize\rmfamily\bfseries\color{black}}
  {\thesubsection}{0.5em}{}

\titleformat{\subsubsection}
  {\small\sffamily\bfseries\color{black}}
  {\thesubsubsection}{0.5em}{}

\titlespacing*{\section}{0pt}{1.0em}{0.4em}
\titlespacing*{\subsection}{0pt}{0.7em}{0.25em}
\titlespacing*{\subsubsection}{0pt}{0.5em}{0.15em}

% ============================================================
%  CATEGORY TAG — coloured pill label
% ============================================================
\newcommand{\cattag}[2]{%
  \tikz[baseline=(tag.base)]{%
    \node[fill=#1, text=white, rounded corners=2.5pt,
          inner xsep=5pt, inner ysep=1.5pt,
          font=\footnotesize\sffamily\bfseries] (tag) {#2};}}

% ============================================================
%  LOGO HELPER — uses actual file if found, else coloured badge
%  Usage: \logoimg{jlu}{jluDeepBlue}{吉林大学}
%  It searches these paths in order:
%    logo/<name>.pdf, logo/<name>.png,
%    logos/<name>.pdf, logos/<name>.png,
%    <name>.pdf, <name>.png
% ============================================================
\newcommand{\logobadge}[2]{%
  \tikz[baseline=-0.5ex]{%
    \node[fill=#1, text=white, rounded corners=3pt,
          inner xsep=6pt, inner ysep=3pt,
          font=\small\sffamily\bfseries]{#2};}}

\newcommand{\logoimg}[4][1.9em]{% [height]{stem}{color}{fallback-text}
  \IfFileExists{logo/#2.pdf}%
    {\includegraphics[height=#1]{logo/#2}}%
    {\IfFileExists{logo/#2.png}%
      {\includegraphics[height=#1]{logo/#2}}%
      {\IfFileExists{logos/#2.pdf}%
        {\includegraphics[height=#1]{logos/#2}}%
        {\IfFileExists{logos/#2.png}%
          {\includegraphics[height=#1]{logos/#2}}%
          {\IfFileExists{#2.pdf}%
            {\includegraphics[height=#1]{#2}}%
            {\IfFileExists{#2.png}%
              {\includegraphics[height=#1]{#2}}%
              {\logobadge{#3}{#4}}}}}}}}

% ============================================================
%  HIGHLIGHT BOX — compact
% ============================================================
\newtcolorbox{highlightbox}[1][]{%
  enhanced, breakable,
  colback=jluLightBg,
  colframe=jluDeepBlue,
  coltitle=white,
  fonttitle=\small\sffamily\bfseries,
  boxrule=0.6pt,
  arc=3pt,
  left=5pt, right=5pt, top=4pt, bottom=4pt,
  title=#1}

% ============================================================
%  DOCUMENT
% ============================================================
\begin{document}
\thispagestyle{fancy}

% ======================  TITLE BLOCK (full width) ============
% --- title ---
\begin{center}
  {\fontsize{20}{24}\selectfont\rmfamily\bfseries\color{black}%
    Learning to Stylize by Learning to Destylize:\par}
  \vspace{1pt}
  {\fontsize{16}{20}\selectfont\rmfamily\bfseries\color{black}%
    A Scalable Paradigm for Supervised Style Transfer\par}
  \vspace{6pt}

  % --- authors ---
  {\normalsize\rmfamily
    \textbf{Ye Wang}\textsuperscript{1}\quad
    \textbf{Zili Yi}\textsuperscript{2}\quad
    \textbf{Yibo Zhang}\textsuperscript{1,3}\quad
    \textbf{Peng Zheng}\textsuperscript{1,3}\quad
    \textbf{Xuping Xie}\textsuperscript{1}\quad
    \textbf{Jiang Lin}\textsuperscript{2}\\[2pt]
    \textbf{Yijun Li}\textsuperscript{4}\quad
    \textbf{Yilin Wang}\textsuperscript{4,$\star\dagger$}\quad
    \textbf{Rui Ma}\textsuperscript{1,5,$\star$}}\\[4pt]
  % --- affiliations ---
  {\footnotesize\rmfamily\color{black}
    \textsuperscript{1}Jilin University\quad
    \textsuperscript{2}Nanjing University\quad
    \textsuperscript{3}Shanghai Innovation Institute\\[1pt]
    \textsuperscript{4}Adobe\quad
    \textsuperscript{5}Engineering Research Center of Knowledge-Driven Human-Machine Intelligence, MOE, China}\\[2pt]
  {\footnotesize\rmfamily\color{jluGray}
    $\star$ Corresponding authors \quad\quad $\dagger$ Project lead}
\end{center}

\vspace{-1em}
\noindent{\color{jluDeepBlue}\rule{\textwidth}{0.7pt}}
\vspace{-1em}

% --- abstract (full width) ---
\begin{center}
\begin{tcolorbox}[
  enhanced,
  width=0.92\textwidth,
  colback=jluDeepBlue!6!white,
  colframe=jluDeepBlue!35!white,
  boxrule=0.6pt,
  arc=3pt,
  left=6pt,right=6pt,top=5pt,bottom=5pt
]
  {\small
  \noindent{\sffamily\bfseries\color{jluDeepBlue}Abstract\quad}%
  Write a concise abstract here (150--250 words). Briefly state the
  problem, your approach, key results, and the main conclusion. Avoid
  citations and undefined abbreviations in the abstract. This template
  uses a compact single-column layout with tight margins inspired by
  top-tier ML papers, allowing clear presentation of text, figures,
  and equations in one continuous flow.}
\end{tcolorbox}
\end{center}

\vspace{0.4em}
\begin{center}
  {\footnotesize\sffamily\bfseries\color{jluDeepBlue}Keywords:\;}
  {\footnotesize keyword one, keyword two, keyword three, keyword four}
  \\[4pt]
  \cattag{jluDeepBlue}{Topic A}\;\;
  \cattag{jluRed}{Topic B}\;\;
  \cattag{jluTeal}{Topic C}\;\;
  \cattag{jluAmber}{Topic D}
\end{center}

\vspace{0.8em}

% ====================== SINGLE-COLUMN BODY ======================

% ============================================================
\section{Introduction}
\label{sec:intro}
% ============================================================

Introduce the research problem and its significance. Provide background
context and clearly state the contributions of your work. The compact
single-column layout keeps readability high while still allowing dense
technical content.

\begin{highlightbox}[Main Contributions]
\begin{enumerate}
  \item \textbf{Contribution 1:} Brief description of the first contribution.
  \item \textbf{Contribution 2:} Brief description of the second contribution.
  \item \textbf{Contribution 3:} Brief description of the third contribution.
\end{enumerate}
\end{highlightbox}

% ============================================================
\section{Related Work}
\label{sec:related}
% ============================================================

Review relevant prior work. Group papers thematically and explain how
your work differs from or builds upon existing methods.

\subsection{Sub-topic A}
Discussion of sub-topic A literature. With the compact layout, related
work sections can accommodate thorough literature review without
consuming excessive page real estate.

\subsection{Sub-topic B}
Discussion of sub-topic B literature and connections to your approach.

% ============================================================
\section{Methods}
\label{sec:methods}
% ============================================================

Describe your proposed method in detail.

\subsection{Problem Formulation}
\label{sec:formulation}

Define notation and formally state the problem:
\begin{equation}
  \mathcal{L}(\theta) \!=\! \frac{1}{N}\!\sum_{i=1}^{N}
      \ell\bigl(f_\theta(x_i), y_i\bigr)
    + \lambda\|\theta\|_2^2
  \label{eq:loss}
\end{equation}
where $f_\theta$ is the model parameterised by $\theta$, $(x_i,y_i)$
are training samples, and $\lambda$ is the regularisation coefficient.

\subsection{Model Architecture}
\label{sec:architecture}

Explain the architecture or algorithm design. You can use category tags
inline to label components: \cattag{jluTeal}{Encoder}
\cattag{jluRed}{Decoder} \cattag{jluAmber}{Loss}.

\subsection{Training Procedure}
\label{sec:training}

Describe training details, hyper-parameter choices, and optimisation
strategy. Example algorithm block:

\begin{algorithm}[H]
\caption{Training procedure}\label{alg:train}
\begin{algorithmic}[1]
\small
\REQUIRE Dataset $\mathcal{D}$, learning rate $\eta$, epochs $T$
\FOR{$t = 1$ \TO $T$}
  \FOR{mini-batch $(x,y) \sim \mathcal{D}$}
    \STATE $g \leftarrow \nabla_\theta \mathcal{L}(\theta; x, y)$
    \STATE $\theta \leftarrow \theta - \eta \cdot g$
  \ENDFOR
\ENDFOR
\end{algorithmic}
\end{algorithm}

% ============================================================
\section{Experiments}
\label{sec:experiments}
% ============================================================

\subsection{Experimental Setup}
\label{sec:setup}

Describe datasets, evaluation metrics, baselines, and hardware/software
environment. The compact format works well for dense experimental
sections with multiple tables and figures.

\subsection{Main Results}
\label{sec:results}

\begin{table}[H]
  \centering
  \caption{Comparison on Benchmark X.}
  \label{tab:main_results}
  \small
  \rowcolors{2}{jluLightBg}{white}
  \begin{tabular}{@{}lccc@{}}
    \toprule
    \textbf{\sffamily Method}
      & \textbf{\sffamily Acc.\,(\%)}
      & \textbf{\sffamily F1}
      & \textbf{\sffamily \#P\,(M)} \\
    \midrule
    Baseline A    & 85.3 & 0.841 & 12.5 \\
    Baseline B    & 87.1 & 0.862 & 24.0 \\
    Baseline C    & 88.0 & 0.873 & 30.1 \\
    \rowcolor{blue!8}
    \textbf{Ours} & \textbf{\color{jluRed}89.7}
                  & \textbf{\color{jluRed}0.891} & 18.2 \\
    \bottomrule
  \end{tabular}
\end{table}

\subsection{Ablation Study}
\label{sec:ablation}

\begin{table}[H]
  \centering
  \caption{Ablation study on key components.}
  \label{tab:ablation}
  \small
  \begin{tabular}{@{}lcc@{}}
    \toprule
    \textbf{\sffamily Variant} & \textbf{\sffamily Acc.\,(\%)} & \textbf{\sffamily $\Delta$} \\
    \midrule
    Full model           & \textbf{89.7} & ---    \\
    w/o Component A      & 87.2 & \color{jluRed}$-$2.5 \\
    w/o Component B      & 86.5 & \color{jluRed}$-$3.2 \\
    w/o Component C      & 88.4 & \color{jluRed}$-$1.3 \\
    \bottomrule
  \end{tabular}
\end{table}

\begin{highlightbox}[Key Finding]
Removing component B causes the largest drop
(\textbf{\color{jluRed}$-$3.2\%}), confirming its critical role.
\end{highlightbox}

\subsection{Qualitative Analysis}
\label{sec:qualitative}

\begin{figure}[H]
  \centering
  % \includegraphics[width=\linewidth]{figures/qualitative.pdf}
  \fbox{\parbox{0.92\linewidth}{\centering\vspace{1.5em}
    {\small\sffamily\color{jluGray}[Your figure here]}
  \vspace{1.5em}}}
  \caption{Qualitative comparison. Best viewed in colour.}
  \label{fig:qualitative}
\end{figure}

% ============================================================
\section{Discussion}
\label{sec:discussion}
% ============================================================

Interpret results, discuss limitations, and suggest future directions.

\subsubsection{Limitations}
Describe limitations honestly.

\subsubsection{Future Work}
Outline promising follow-up avenues.

% ============================================================
\section{Conclusion}
\label{sec:conclusion}
% ============================================================

Summarise findings and contributions. Restate significance and impact.

% ============================================================
\section*{Acknowledgements}
% ============================================================
{\small This work was supported by [funding agency / grant number].
The authors thank [collaborators] for helpful discussions.}

% ============================================================
%  References
% ============================================================
{\small
% \bibliography{references}
\begin{thebibliography}{9}
\bibitem{ref1}
  A.~Author and B.~Author.
  Title of the first reference.
  \emph{Journal Name}, 1(2):34--56, 2024.

\bibitem{ref2}
  C.~Author, D.~Author, and E.~Author.
  Title of the second reference.
  In \emph{Proc.\ Conference}, pp.\ 78--90, 2025.
\end{thebibliography}
}

% ====================== APPENDIX ======================
\newpage
\appendix
\section{Additional Experimental Details}
\label{app:details}

Supplementary material: extended tables, proofs, or extra figures.

\section{Hyper-parameter Sensitivity}
\label{app:hyperparams}

Additional hyper-parameter analysis or sensitivity plots.

\end{document}